\documentclass[11pt]{article}
\usepackage[margin=1in]{geometry}
\usepackage[T1]{fontenc}
\usepackage[utf8]{inputenc}
\usepackage{amsmath,amssymb,amsthm}
\usepackage{enumitem}

\title{ICPC World Finals 2025\\A. A-Skew-ed Reasoning}
\author{}
\date{}

\begin{document}
\maketitle

\section*{Problem Summary}

We are given a rooted heap-ordered binary tree on the keys $1,2,\dots,n$, where key $i$ is the node
label. We must decide whether this tree can be produced by the skew-heap insertion rule from some
permutation of the keys. If it can, we must output both the lexicographically smallest and the
lexicographically largest such permutations.

\section*{Key Observations}

\begin{itemize}[leftmargin=*]
    \item Because the labels are exactly $1,\dots,n$ and the tree is a min-heap, every subtree root is the
    smallest key in that subtree. So the root of a subtree is the first key of that subtree that can appear
    in the insertion order.
    \item Fix a node $u$. When a later insertion $x>u$ reaches $u$, the skew-heap rule swaps the two
    children of $u$ and then continues recursively into the new left child. Therefore the recursive target
    alternates between the two children of $u$.
    \item Let $A$ be the child subtree that already existed when $u$ itself was inserted, and let $B$ be the
    other child subtree. Then the insertion order inside the subtree of $u$ has only one possible shape:
    some prefix of the order of $B$, then $u$, then the orders of $A$ and the remaining part of $B$
    interleaved in alternating positions.
    \item This immediately gives the feasibility conditions.
    If $A$ is the final left child, then $1 \le |L| \le |R|+1$.
    If $A$ is the final right child, then $|L| \ge |R|$.
    If neither condition holds, the whole tree is impossible.
    \item The root key $u$ is the smallest key in its subtree, so once an orientation is fixed, the
    lexicographic order is decided first by the position of $u$ inside that subtree. Choosing the earliest
    feasible position gives the lexicographically smallest answer; choosing the latest feasible position
    gives the largest answer.
\end{itemize}

\section*{Algorithm}

\begin{enumerate}[leftmargin=*]
    \item Compute every subtree size bottom-up. This is possible because every child label is larger than its
    parent label.
    \item For each node, test the two feasible interpretations described above.
    If both are feasible, compare the resulting position of the root inside the local order and record the
    choice for the lexicographically smallest and largest answers separately.
    \item Reconstruct the permutation iteratively.
    For a subtree we do not store a contiguous segment of answer positions; instead we store an arithmetic
    progression of positions.
    The interleaving rule says exactly which positions belong to the prefix of $B$, which single position
    contains the root, which alternating positions belong to $A$, and which alternating positions belong
    to the remaining part of $B$.
    \item Applying this split recursively yields the entire minimum permutation, and repeating it with the
    other orientation choices yields the maximum permutation.
\end{enumerate}

\section*{Correctness Proof}

We prove that the algorithm returns the correct answer.

\paragraph{Lemma 1.}
For any valid subtree rooted at $u$, after fixing which child was the previously existing heap at the
moment $u$ was inserted, the insertion order inside that subtree has the interleaving pattern described in
the observations.

\paragraph{Proof.}
The key $u$ is inserted when it becomes the new minimum of its subtree, so at that moment the old heap
becomes one child of $u$ and the other child is empty. Every later key in the same subtree is larger than
$u$, so when it reaches $u$ the insertion rule swaps the children and then descends into the new left
child. Therefore the recursive destination alternates between the two child subtrees. Before the two
streams can alternate, the larger one may need an initial prefix so that both streams still have elements
available afterwards. This is exactly the stated pattern. \qed

\paragraph{Lemma 2.}
The size inequalities used by the algorithm are necessary and sufficient for a node to be feasible.

\paragraph{Proof.}
If $A$ is the final left child and $B$ is the final right child, then after the root there must be enough
elements of $B$ to occupy the alternating positions around the elements of $A$, so $|A|$ can exceed
$|B|$ by at most $1$, and $A$ cannot be empty. This gives $1 \le |L| \le |R|+1$.
The case where $A$ is the final right child is symmetric and yields $|L| \ge |R|$.

Conversely, if one of these inequalities holds, the arithmetic-progression split produced by the
algorithm assigns a valid set of answer positions to each of the two child streams and to the root.
Recursing on the children therefore produces a valid insertion order for the whole subtree. \qed

\paragraph{Theorem.}
The algorithm outputs the correct answer for every valid input.

\paragraph{Proof.}
By Lemma 2, the algorithm rejects exactly the impossible nodes, so it prints \texttt{impossible} exactly
when no valid insertion order exists.

Assume the tree is feasible. By Lemma 1, every valid insertion order of a subtree must use one of the
two local orientations considered by the algorithm, and by Lemma 2 the recursive placement generated
from that orientation is valid. Therefore both reconstructed permutations are valid.

For lexicographic minimality, note that the subtree root is the smallest key in that subtree. Hence among
all valid local orientations, the one that places the root earlier makes the whole subtree lexicographically
smaller, regardless of what happens later. The algorithm chooses exactly that orientation for the minimum
answer, and the symmetric later-root choice for the maximum answer. Applying this argument
independently at every subtree proves that the produced permutations are the lexicographically smallest
and largest valid ones. \qed

\section*{Complexity Analysis}

Computing subtree sizes is $O(n)$. The feasibility pass is also $O(n)$. The reconstruction touches every
node a constant number of times, so it is $O(n)$. The memory usage is $O(n)$.

\section*{Implementation Notes}

\begin{itemize}[leftmargin=*]
    \item The reconstruction is written iteratively to avoid recursion depth problems at $n=2\cdot 10^5$.
    \item Each stack frame stores an arithmetic progression of answer indices rather than an explicit list of
    positions, which keeps the reconstruction linear.
\end{itemize}

\end{document}
